\chapter{Thiết kế và cài đặt hệ thống}
\label{chap:design}

\section{Kiến trúc tổng thể}

Hệ thống được tổ chức thành một chuỗi xử lý tuần tự cho từng ảnh hoặc từng khung hình video
(Hình~\ref{fig:pipeline}). Trên video, thư viện Ultralytics kết hợp YOLO11-seg với ByteTrack để trả về hộp, lớp, độ tin cậy, ID và mặt nạ của từng xe. Phần dựng bao chỉ nhận tập điểm biên, nên có thể dùng với mô hình khác nếu mô hình đó cũng cung cấp mặt nạ đối tượng.

\begin{figure}[H]
\centering
\begin{tikzpicture}[grp/.style={draw, dashed, rounded corners, inner sep=5pt}]
  \node (in)   [blk] at (0,0)      {Ảnh / khung hình video};
  \node (yolo) [bb]  at (0,-1.3)   {YOLO11-seg (large)\\ lọc 4 lớp, NMS};
  \node (det)  [bb]  at (0,-2.6)   {Hộp, lớp, độ tin cậy\\ + mặt nạ từng đối tượng};
  % Chế độ video
  \node (trk)  [nk]  at (-3.2,-4.5) {ByteTrack có sẵn:\\ chỉ giữ xe có ID};
  \node (pv)   [nk]  at (-3.2,-5.9) {Đường biên mặt nạ\\ $\rightarrow$ tập điểm $P$\\ (thiếu mặt nạ: vẽ hộp)};
  \node (hv)   [nk]  at (-3.2,-7.4) {Dựng bao bằng\\ hai bản cài đặt};
  \node (ov)   [hd]  at (-3.2,-8.8) {Video có ID + bao;\\ CSV đặc trưng};
  % Chế độ đánh giá
  \node (pe)   [nk]  at (3.2,-4.5)  {Mọi dự đoán trên ảnh có nhãn\\ (không theo vết)};
  \node (he)   [nk]  at (3.2,-5.9)  {Đường biên mặt nạ $\rightarrow$ bao\\ trực giao, bao lồi thường};
  \node (ge)   [nk]  at (3.2,-7.4)  {So từng đầu ra\\ với nhãn chuẩn};
  \node (oe)   [hd]  at (3.2,-8.8)  {Số liệu, biểu đồ,\\ ảnh đối chiếu};
  \draw [arrow] (in) -- (yolo);
  \draw [arrow] (yolo) -- (det);
  \draw [arrow] (det.south) -- (trk.north);
  \draw [arrow] (det.south) -- (pe.north);
  \draw [arrow] (trk) -- (pv);
  \draw [arrow] (pv) -- (hv);
  \draw [arrow] (hv) -- (ov);
  \draw [arrow] (pe) -- (he);
  \draw [arrow] (he) -- (ge);
  \draw [arrow] (ge) -- (oe);
  \begin{scope}[on background layer]
    \node [grp, fit=(trk)(ov), label={[font=\small\bfseries]below:Chế độ video}] {};
    \node [grp, fit=(pe)(oe), label={[font=\small\bfseries]below:Chế độ đánh giá}] {};
  \end{scope}
\end{tikzpicture}
\caption{Sơ đồ hai chế độ chạy của hệ thống}
\label{fig:pipeline}
\end{figure}

Hệ thống có hai chế độ chạy:
\begin{enumerate}
  \item \textbf{Chế độ video}: gọi \texttt{model.track()} để chạy YOLO11-seg cùng ByteTrack, rồi chỉ dựng bao cho các đối tượng
        đã có ID. Đối tượng thiếu mặt nạ được vẽ bằng hộp. Kết quả được vẽ ra hai video riêng và đặc trưng hình học của từng đối tượng được ghi ra CSV.
  \item \textbf{Chế độ đánh giá}: chạy trên tập ảnh có nhãn polygon chuẩn, không qua bộ theo vết. Sáu loại đầu ra (hộp, vùng hình chữ nhật, bao lồi thường, mặt nạ YOLO, hai bản cài đặt bao trực giao)
        được so với nhãn để xuất số liệu, biểu đồ, ảnh đối chiếu.
\end{enumerate}
Do chế độ video chỉ dựng bao cho các đối tượng có ID, số liệu trên ảnh không xác nhận trực tiếp chất lượng của toàn bộ đầu ra video.

\section{Hai bản cài đặt dựng bao}
\label{sec:impl_hull}

Hai bản cài đặt đều làm việc với tối đa tám điểm mốc và bốn chuỗi bậc thang ở mục~\ref{sec:staircase}.
Đây là mã của nhóm, không phải mã của các tác giả bài báo.

\subsection{So sánh hai bản cài đặt}

Bảng~\ref{tab:algo_compare} so sánh hai bản cài đặt của nhóm. Bản QuickHull sắp xếp lại tập điểm ở mỗi lần gọi và có thể cần nhiều lần gọi khi số điểm trên biên lớn, nên có cận xấu nhất $O(n^2\log n)$; bản Graham có cận $O(n\log n)$.



\begin{table}[H]
\centering
\caption{So sánh hai bản cài đặt của nhóm}
\label{tab:algo_compare}
\begin{tabular}{>{\raggedright\arraybackslash}p{4.2cm} >{\raggedright\arraybackslash}p{5.2cm} >{\raggedright\arraybackslash}p{5.2cm}}
\toprule
\textbf{Tiêu chí} & \textbf{Bản QuickHull} & \textbf{Bản Graham} \\
\midrule
Chiến lược & Đệ quy một nhánh & Sắp xếp rồi quét một lượt \\
Độ phức tạp xấu nhất & $O(n^2\log n)$ & $O(n\log n)$ \\
Bộ nhớ phụ xấu nhất & $O(n^2)$ (danh sách điểm ở mỗi tầng đệ quy) & $O(n)$ \\
Kết quả trên dữ liệu thử & \multicolumn{2}{c}{Cùng một vùng bao trên dữ liệu thử} \\
\bottomrule
\end{tabular}
\end{table}

Hai bản cài đặt chưa bảo đảm bao tối tiểu trên mọi đầu vào; chẳng hạn, với hai điểm chéo nhau, mã trả về hình chữ nhật thay vì đường bậc thang.

\section{Chi tiết cài đặt}

Hệ thống được viết bằng Python. Để hai bản cài đặt được so sánh công bằng, cả hai đều được gọi qua cùng một hàm chuẩn hoá:
\begin{enumerate}
  \item Ép toạ độ về số nguyên và loại các điểm trùng nhau.
  \item Nếu còn ít hơn 2 điểm phân biệt, trả về chính các điểm đó. Với 2 hoặc 3 điểm, hai bản cài đặt dùng chung nhánh O-Graham để mọi cạnh vẫn song song với trục.
  \item Gọi bản QuickHull hoặc bản Graham.
  \item Thêm lại điểm đầu vào cuối dãy đỉnh để khép kín đa giác.
\end{enumerate}
Sau đó, hệ thống tính diện tích, chu vi, tỉ lệ lấp đầy so với hộp và số đỉnh của bao.
Bộ đếm thời gian chỉ bao quanh bước dựng bao, không tính inference YOLO, vẽ hay ghi tệp.
Trên video, ID được lấy từ \texttt{result.boxes.id}; hộp, lớp, độ tin cậy và mặt nạ được đọc theo cùng chỉ số trong kết quả do Ultralytics trả về. Nhờ đó, bao được dựng từ mặt nạ tương ứng với ID đang hiển thị.

Nhãn chuẩn ở định dạng YOLO polygon (đa giác có toạ độ được chia cho chiều rộng và chiều cao ảnh để đưa về khoảng từ 0 đến 1): mỗi polygon được nhân với kích thước ảnh, làm tròn rồi tô thành mặt nạ tại độ phân giải gốc; ảnh có nhãn phương tiện chỉ ở dạng hộp bị loại.
